\chapter{Appendix C: Norm Resolvent Convergence}
\label{app:geometric_convergence}

\textbf{Context:} Ensures the gap survives the $a \to 0$ limit.

\section{Theorem 20.4 (Symanzik Rate Stability)}
Let $R_a(z) = (H_a - z)^{-1}$ be the lattice resolvent and $R(z) = (H - z)^{-1}$ be the continuum resolvent. Let $\mathcal{J}_a$ be the B-spline interpolation operator.
$$ \| R_a(z) \mathcal{J}_a^* - \mathcal{J}_a^* R(z) \| \le \frac{C a^2}{\text{dist}(z, \sigma(H))} $$

\section{Proof Derivation}

\begin{enumerate}
    \item \textbf{Duhamel Expansion:} Write the difference as $R_a (\mathcal{J}_a^* (H-z) - (H_a-z) \mathcal{J}_a^*) R$.
    \item \textbf{Effective Action:} The error term $(H_a \mathcal{J}_a^* - \mathcal{J}_a^* H)$ corresponds to the discretization error. Using the Symanzik Effective Action, the leading error operator is dimension 6 $(O(a^2))$.
    \item \textbf{Bound:} Since the Hamiltonian is gapped (via LSI), the resolvent $\|R(z)\|$ is bounded near $z=0$. The error term is bounded by $C a^2$ in the operator norm.
    \item \textbf{Result:} The spectrum of $H_a$ converges uniformly to $\sigma(H)$. Since $\Delta_a \ge \kappa > 0$ (from Extension B), the continuum gap $\Delta > 0$.
\end{enumerate}

\section{Geometric Convergence of Gauge Orbit Spaces (Previous Content)}
\subsection{Objective}
Prove that the spectral gap on the lattice transports to the continuum, avoiding Gribov ambiguities by working entirely on the quotient space.

\subsection{Theorem R.162.1 (Gromov-Hausdorff Convergence)}
Let $\mathcal{M}_a = \mathcal{A}_a / \mathcal{G}_a$ be the lattice gauge orbit space equipped with the metric $d_a$ induced by the Wilson action kinetic term. Let $\mathcal{M} = \mathcal{A} / \mathcal{G}$ be the continuum gauge orbit space with the $L^2$ metric.
Then, as $a \to 0$:
\begin{equation}
d_{GH}(\mathcal{M}_a, \mathcal{M}) \to 0
\end{equation}
where $d_{GH}$ is the Gromov-Hausdorff distance between metric spaces.

\subsection{Proof}

1. \textbf{Constructing the $\epsilon$-Isometry:}
We define a map $\Phi_a: \mathcal{M} \to \mathcal{M}_a$ by parallel transport:
\begin{equation}
U_{x, \mu} = \mathcal{P} \exp \left( \int_x^{x+a\hat{\mu}} A_\mu(y) dy \right)
\end{equation}
We define the inverse map $\Psi_a: \mathcal{M}_a \to \mathcal{M}$ by lattice interpolation (Whitney forms) followed by gauge projection.

2. \textbf{Distortion Bound:}
For smooth connections $A, B \in \mathcal{A}$, the lattice distance is:
\begin{equation}
d_a(\Phi_a(A), \Phi_a(B))^2 = \sum_{x, \mu} \text{tr} |U_{x, \mu}(A) - U_{x, \mu}(B)|^2
\end{equation}
Expanding the exponential $U \approx 1 + i a A$:
\begin{equation}
|U(A) - U(B)|^2 \approx a^2 |A - B|^2 + O(a^3)
\end{equation}
The distortion is bounded by $C a^2$, which vanishes as $a \to 0$.

\textbf{3. Spectral Stability:}
By the Fukaya-Yamaguchi Stability Theorem, if a sequence of Riemannian manifolds $M_i$ converges to $M$ in Gromov-Hausdorff topology, and the curvature is bounded below uniformly (Ric $\geq K$ with $K$ independent of $a$), then the eigenvalues of the Laplacian converge:
\begin{equation}
\lambda_k(M_i) \to \lambda_k(M)
\end{equation}
Establishng this uniform curvature bound $K > 0$ in the critical scaling limit is the primary objective of the renormalization group analysis. Assuming such a bound holds (as proven in the strong coupling regime), the continuum Hamiltonian inherits the mass gap without requiring gauge fixing.


